\chapter{Optimization, Generalization, and Scaling Intuition}
\label{sec:ch2-optgen}

\section*{Setting the Scene}
\addcontentsline{toc}{section}{Setting the Scene}

\paragraph{The problem.} Chapter~1 told us \emph{what} we want a
language model to do: assign a probability distribution to the next
token. It said almost nothing about how to find the parameters that
realise such a model. We now face that question. The model has, on
the small side, a hundred million parameters; on the large side, a
trillion. The loss surface is a function of all of them at once.
There is no closed-form solution. There is no convex landscape to
descend. There is, instead, a vast and largely unmapped territory,
and our only tool for navigating it is the gradient at our feet.

\paragraph{How we got here.} The mathematics of descending a
function by following its gradient is older than electricity. In
1847 the French mathematician Augustin-Louis Cauchy described the
method of steepest descent in a short note to the
\emph{Comptes Rendus}. Cauchy's problem was small enough that the
gradient could be evaluated exactly. Ours is not. The trick that
makes our problem tractable is a much later one: replace the true
gradient with a noisy estimate computed from a small batch of
examples. Robbins and Monro, in 1951, gave the conditions under
which such a noisy iteration converges. They called it stochastic
approximation. We call it stochastic gradient descent.

The other trick we needed was a way to compute the gradient of a
deep nonlinear function at all. The chain rule of calculus answers
this in principle. The implementation is called
back-propagation. Paul Werbos, in his 1974 Harvard PhD thesis,
wrote down a recognisable version of it. Rumelhart, Hinton, and
Williams brought it to the field's attention in their 1986 paper
in \emph{Nature}. Without back-propagation a network of more than
a few layers is computationally inert. With it, the gradient
computation is no longer a bottleneck against depth: the cost of
computing the gradient is roughly the cost of computing the forward
pass, so adding layers raises training cost only linearly in the
depth.

For thirty years the field then argued about which optimiser was
best. Momentum, RMSProp, Adagrad, Adadelta. Diederik Kingma and
Jimmy Ba published Adam in 2014, combining momentum and adaptive
per-parameter step sizes. AdamW, the variant in current use, came
from Loshchilov and Hutter in 2017 and fixed a subtle but
consequential issue with how Adam handled weight decay. By 2020
the question of \emph{which} optimiser to use for large-scale
Transformer pretraining was largely settled at AdamW. The
question of \emph{how much to spend on what} took its place. Jared
Kaplan and his coauthors at OpenAI, in 2020, observed that loss
followed a power law in the three quantities of interest:
parameter count, dataset size, and compute. Hoffmann and her
colleagues at DeepMind, in 2022, ran a more careful study and
concluded that previous large models had been trained on too few
tokens for their parameter counts. They called this Chinchilla
scaling. Today's compute-optimal recipes are descendants of that
single argument.

\paragraph{Where we stand.} A modern training recipe is a
sequence of decisions, not an algorithm. Choose AdamW. Pick a
peak learning rate, with a linear warmup from zero over the first
few thousand steps and a cosine decay to a small fraction of peak
over the remaining steps. Clip gradients at a global norm of one.
Train in mixed precision (bf16 for most arithmetic, fp32 for
optimiser state). Use gradient accumulation to fit a desired
effective batch size. Check the loss curve daily and the gradient
norm hourly. Pick a model size and a dataset size at the
Chinchilla-optimal ratio for your compute budget. The mathematics
in the rest of this chapter is the bookkeeping for these choices:
the SGD-as-stochastic-differential-equation argument that
explains why warmup helps, the AdamW update rule with its bias
corrections, the mixed-precision memory budget that determines
what fits on a GPU, the Lagrangian that produces the
compute-optimal model and data sizes from a fixed budget.

\paragraph{What goes wrong.} The dramatic failures of training
are visible. A loss spike that never recovers. A gradient norm
that blows up after step 50,000. A divergence that turns a
beautiful run into a wasted week of cluster time. These look bad,
and they cost money, but they are the easy ones. The engineer
who notices them can usually diagnose them: a learning rate too
high, a batch too small, a precision too low, an outlier example
that wedged itself into the training distribution.

The harder failures are the quiet ones. A model that trains to a
respectable validation loss and yet generalises poorly to the
distribution it will actually meet at deployment. A model that is
slightly too small for the data it has seen, and so behaves as if
it is memorising rather than learning. A model that is slightly
too large for the data it has seen, and so behaves as if it has
discovered a regularity that is not there. The double-descent
curves we will discuss are not theoretical curiosities; they are
the visible signature of the implicit regularisation done by
overparameterisation, and getting the signature wrong is what
turns a training run from a model into an expensive monument.

There is a more sociological failure mode worth naming early. The
field has acquired a habit of explaining results by appealing to
``flat minima'' and ``benign overfitting'' as if these were
mechanisms rather than observations. They are observations. We
will treat them as such. The math in the rest of the chapter is
not a sufficient theory of generalisation; nobody has one. It is
a working set of identities and budget equations that, taken
together, let an engineer reason about the regimes in which the
training they are running is likely to behave well, and the
regimes in which it is not.

\section{Learning Objectives}

By the end of this chapter, you should be able to:
\begin{itemize}
  \item Derive SGD, momentum, and Adam/AdamW updates from first principles and
        explain the role of bias correction.
  \item Reason quantitatively about learning-rate schedules (linear warmup,
        cosine decay, inverse-square-root) and their engineering tradeoffs.
  \item Describe why gradient clipping is necessary for deep Transformer
        stacks and choose a clipping threshold defensibly.
  \item State the Kaplan and Chinchilla scaling laws and derive the
        compute-optimal model/data ratio.
  \item Explain double descent and benign overfitting, and articulate why
        overparameterization can \emph{improve} generalization in practice.
  \item Evaluate mixed-precision training, loss scaling, gradient
        accumulation, and checkpointing as systems-level mechanisms.
  \item Design a reproducibility protocol (seeds, deterministic kernels,
        data ordering, regression gates) for production training pipelines.
\end{itemize}

\section{LLM Components Covered}

\textbf{Training dynamics}
\begin{itemize}
  \item Stochastic gradient estimators and their noise structure.
  \item First-order adaptive optimizers (Adam, AdamW) and their
        moment-estimate machinery.
  \item Learning-rate schedules, warmup, and decay.
  \item Gradient clipping, gradient accumulation, and loss scaling.
\end{itemize}

\textbf{Generalization machinery}
\begin{itemize}
  \item Overparameterization and implicit regularization.
  \item Double descent and benign overfitting.
  \item Flat-vs-sharp minima in the large-batch regime.
\end{itemize}

\textbf{Systems relevance}
\begin{itemize}
  \item Training instability, divergence, and regressions.
  \item Reproducibility under distributed training.
  \item Scaling-law-driven capacity planning.
  \item Cost / quality tradeoff under fixed compute budgets.
\end{itemize}

\section{Conceptual Core: Optimization Is Not Inference}

Training an LLM is, mechanically, a large-scale numerical optimization
problem. The learner does not \emph{reason} about the target; it performs a
long sequence of gradient updates against an empirical risk. The objective
from Chapter~1, indexed over the $N_{\mathrm{tok}}$ (context, token) pairs in
the training set, is the empirical cross-entropy
\begin{equation}
\label{eq:ch2-empirical-loss}
\mathcal{L}(\theta)
  \;=\; \frac{1}{N_{\mathrm{tok}}} \sum_{i=1}^{N_{\mathrm{tok}}}
        -\log p_{\theta}\!\left( x^{(i)}_{t_i} \,\middle|\, x^{(i)}_{<t_i} \right),
\end{equation}
where $i$ indexes (context, token) pairs (Chapter~1, A1.1b) and the
symbol $N$ is reserved for the model parameter count later in the
chapter. The objective is optimized via variants of stochastic
gradient descent:
\begin{equation}
\label{eq:ch2-sgd}
\theta_{t+1} \;=\; \theta_t - \eta_t\, g_t,
\qquad
g_t \;=\; \frac{1}{B}\sum_{i \in \mathcal{B}_t}
          \nabla_\theta \bigl[-\log p_{\theta}(x^{(i)}_{t_i} \mid x^{(i)}_{<t_i})\bigr],
\end{equation}
where $\eta_t$ is the learning rate and $\mathcal{B}_t$ is a minibatch of
size $B$. The key features that distinguish LLM optimization from the
textbook setting are:
\begin{itemize}
  \item \textbf{Dimensionality}: $\theta \in \mathbb{R}^d$ with
        $d \sim 10^{9}$--$10^{12}$ for frontier models.
  \item \textbf{Non-convexity}: the loss surface has a combinatorial
        number of saddle points, plateaus, and basins.
  \item \textbf{Stochasticity}: under uniform random sampling of the
        minibatch from the training set, $g_t$ is an unbiased estimator
        of $\nabla \mathcal{L}$, with covariance scaling as $\Sigma/B$
        provided the per-example gradients have bounded variance and
        are not heavily correlated within the batch (e.g.\ shuffled
        sampling rather than block sampling).
  \item \textbf{Distributed execution}: gradients are aggregated across
        thousands of accelerators, making single-step cost high and
        divergence expensive.
\end{itemize}
That modern deep networks reliably converge under these conditions is an
empirical observation about the loss landscape, not a theorem.

\section{Notation and Standing Assumptions}
\label{sec:ch2-notation}

\begin{notationbox}
Throughout this chapter the parameter vector is
$\theta \in \R^{d}$ with $d \sim 10^{9}$--$10^{12}$; $\mathcal{L}(\theta)$
is the empirical cross-entropy of Eq.~\eqref{eq:ch2-empirical-loss};
$g_t$ is the stochastic minibatch gradient with batch size $B$ and
conditional covariance $\Sigma(\theta)$; $\eta_t$ is the learning rate
at step $t$; $m_t, v_t$ are first- and second-moment accumulators with
decays $\beta_1, \beta_2 \in (0,1)$ and numerical floor
$\epsilon > 0$; $\lambda \ge 0$ is decoupled weight decay. Schedule
parameters are warmup length $T_w$, total horizon $T_{\max}$, peak
rate $\eta_{\max}$, and floor $\eta_{\min}$. Gradient-clip threshold
is $c > 0$. For scaling laws, $N$ is parameter count, $D$ is training
tokens, $C$ is compute in FLOPs, and $(A, B, E, \alpha, \beta)$ are
the Chinchilla fit constants of Eq.~\eqref{eq:ch2-chinchilla}.
Standing assumptions are:
\textbf{(A2.1)} the SGD-to-SDE approximation
(Eq.~\eqref{eq:ch2-langevin}) holds only when $\eta$ is small
relative to curvature, $\Sigma$ is approximately locally stationary,
and $\xi_t$ is approximately Gaussian by a CLT argument over the
minibatch;
\textbf{(A2.2)} Chinchilla exponents $\alpha, \beta$ were fit on
properly tuned baselines and extrapolate only within the compute range
spanned by the fit (roughly $10^{19}$--$10^{24}$ FLOPs);
\textbf{(A2.3)} the memory budget
(Eq.~\eqref{eq:ch2-mem-budget}) assumes standard mixed-precision
AdamW with per-worker full optimizer state; ZeRO-style sharding is
layered on top in Chapter~7 and changes only the per-worker terms, not
the total.
\end{notationbox}

\section{Stochastic Gradient Descent and Momentum}

\subsection{SGD as a Noisy Dynamical System}

Write the minibatch gradient as the true gradient plus noise,
$g_t = \nabla \mathcal{L}(\theta_t) + \xi_t$, with
$\mathbb{E}[\xi_t]=0$ and $\mathrm{Cov}(\xi_t) \approx \Sigma(\theta_t)/B$.

\Assum\ Under A2.1 the discrete-time recursion of
Eq.~\eqref{eq:ch2-sgd} is approximated by a continuous-time SDE
\emph{only} when three conditions hold simultaneously: (i) $\eta$ is
small relative to the inverse Lipschitz constant of
$\nabla \mathcal{L}$ (so the discretization error per step is
$O(\eta^{2})$); (ii) $\Sigma(\theta_t)$ varies slowly along the
trajectory on the timescale $1/\eta$ (local stationarity of the
diffusion tensor); (iii) $\xi_t$ is approximately Gaussian, justified
by a CLT argument over $B$ i.i.d.\ per-example gradients for moderate
$B$. Outside these conditions — very large steps, sharp curvature,
or heavy-tailed per-example gradients — the SDE picture is
illustrative rather than predictive.

Under these assumptions, the continuous-time limit of
Eq.~\eqref{eq:ch2-sgd} is the Langevin SDE
\begin{equation}
\label{eq:ch2-langevin}
d\theta_t \;=\; -\nabla \mathcal{L}(\theta_t)\, dt
                + \sqrt{\eta\,\Sigma(\theta_t)/B}\; dW_t,
\end{equation}
which reveals two systems-relevant facts:
(i)~the stationary distribution is \emph{not} concentrated on the loss
minimum; it is biased toward flatter regions, which correlates with better
generalization~\citep{keskar2017large};
(ii)~the effective ``noise temperature'' scales with $\eta/B$, so doubling
batch size roughly halves gradient noise — this is the basis for the
linear scaling rule in large-batch training~\citep{goyal2017accurate}.

\Emp\ The ``flat minima generalize better'' connection is an
empirical correlation supported by large-batch experiments on image
and language workloads; it is not a theorem and admits known
counterexamples under reparameterization~\citep{dinh2017sharp}. Treat
it as a design-time heuristic, not a guarantee.

\subsection{Heavy-Ball and Nesterov Momentum}

Momentum introduces an exponentially-weighted moving average of past
gradients:
\begin{align}
m_t &= \beta\, m_{t-1} + (1-\beta)\, g_t, \label{eq:ch2-momentum-m} \\
\theta_{t+1} &= \theta_t - \eta_t\, m_t.
\end{align}
Geometrically, momentum acts as a low-pass filter on the gradient signal:
high-frequency noise is attenuated while persistent descent directions
accumulate. Typical values are $\beta \in [0.9, 0.99]$. Nesterov's
look-ahead variant evaluates the gradient at the extrapolated point
$\theta_t - \eta \beta m_{t-1}$, which admits a stronger convergence
rate on convex problems but offers modest gains on LLM objectives.

\section{Adam and AdamW}

\subsection{Why Adaptive Per-Parameter Rates}

In a Transformer, embedding vectors for frequent tokens see large,
frequent gradients while embeddings for rare tokens see small, infrequent
ones. A single global $\eta$ cannot serve both regimes well: it is
either too aggressive for frequent tokens (overshooting) or too
conservative for rare ones (starvation). Adaptive methods scale the
update for each parameter by an online estimate of its gradient
magnitude~\citep{kingma2015adam}.

\subsection{Adam Update Rule}

Let $g_t = \nabla_\theta \mathcal{L}(\theta_t)$ be the minibatch
gradient. Adam maintains two exponential moving averages:
\begin{align}
m_t &= \beta_1\, m_{t-1} + (1-\beta_1)\, g_t
       & \text{(first moment)}, \label{eq:ch2-adam-m}\\
v_t &= \beta_2\, v_{t-1} + (1-\beta_2)\, g_t^2
       & \text{(second moment)},
\end{align}
where $g_t^2$ is elementwise. Because $m_0 = v_0 = 0$, these estimates
are biased toward zero in the first several steps. Adam corrects this
bias analytically.

\begin{derivation}[for Adam bias correction]
Unroll Eq.~\eqref{eq:ch2-adam-m} from $m_0 = 0$:
\begin{align*}
m_t &= (1-\beta_1) \sum_{k=1}^{t} \beta_1^{\,t-k}\, g_k.
\end{align*}
Take expectations term by term, writing $\mu_k := \E[g_k]$:
\begin{align*}
\E[m_t]
  &= (1-\beta_1) \sum_{k=1}^{t} \beta_1^{\,t-k}\,\mu_k.
\end{align*}
\textbf{Assume (local stationarity)} that $\mu_k$ varies slowly on the
timescale $1/(1-\beta_1)$, so $\mu_k \approx \mu_t$ for all $k$ whose
weight $\beta_1^{t-k}$ is non-negligible. Pulling $\mu_t$ out of the
sum and using the geometric series
$\sum_{k=1}^{t}\beta_1^{\,t-k} = (1-\beta_1^{\,t})/(1-\beta_1)$ gives
\begin{align*}
\E[m_t] \;\approx\; (1-\beta_1^{\,t})\,\mu_t,
\qquad
\text{so}\qquad
\hat m_t \;:=\; \frac{m_t}{1-\beta_1^{\,t}} \;\text{ is approximately unbiased}.
\end{align*}
The residual drift term $\E[m_t] - (1-\beta_1^{\,t})\mu_t
= (1-\beta_1)\sum_k \beta_1^{t-k}(\mu_k - \mu_t)$ is
$O\!\left((1-\beta_1)\,\|\nabla^{2}\!\mathcal{L}\|\,\eta\,/(1-\beta_1)\right)
= O(\eta\,\|\nabla^{2}\!\mathcal{L}\|)$ under a one-step Taylor
expansion of $\mu_k$, which is small in the warmup regime that motivates
bias correction in the first place. The identical argument applied to
$v_t$ yields $\hat v_t = v_t/(1-\beta_2^{\,t})$.
\end{derivation}

The parameter update becomes
\begin{equation}
\label{eq:ch2-adam-update}
\theta_{t+1} \;=\; \theta_t - \eta_t\,\frac{\hat m_t}{\sqrt{\hat v_t} + \epsilon},
\qquad
\hat m_t = \frac{m_t}{1-\beta_1^{\,t}},
\quad
\hat v_t = \frac{v_t}{1-\beta_2^{\,t}},
\end{equation}
with default hyperparameters $\beta_1 = 0.9$, $\beta_2 = 0.999$,
$\epsilon = 10^{-8}$. For large Transformer training, practitioners
commonly reduce $\beta_2$ to $0.95$ to make the second-moment estimator
more responsive to nonstationarity.

\subsection{AdamW: Decoupled Weight Decay}

Classical L2 regularization adds $\tfrac{\lambda}{2}\|\theta\|^2$ to the
loss, which introduces a $\lambda \theta$ term into $g_t$. In Adam this
term is then divided by $\sqrt{\hat v_t}$, so parameters with small
gradients see disproportionately strong decay — the effective
regularization is entangled with the gradient magnitude. AdamW
\emph{decouples} weight decay by applying it directly to the parameters,
not through the loss gradient~\citep{loshchilov2019decoupled}:
\begin{equation}
\label{eq:ch2-adamw-update}
\theta_{t+1} \;=\; \theta_t
  - \eta_t\!\left(\frac{\hat m_t}{\sqrt{\hat v_t} + \epsilon} + \lambda\,\theta_t\right).
\end{equation}
This subtle change yields consistently better generalization on
Transformer workloads and is the default optimizer for essentially all
production LLM training.

\begin{center}
\begin{tabular}{|l|l|l|}
\hline
Optimizer & Strengths & Risks / Caveats \\
\hline
SGD        & Simple, flatness-biased & Slow without tuning \\
SGD+Mom.   & Accelerates narrow valleys & Still sensitive to $\eta$ \\
Adam       & Fast early progress & Sharp minima, weight-decay bug \\
AdamW      & Decoupled decay, stable & Hyperparameter-sensitive, memory-heavy \\
Lion/Shampoo & Less memory, competitive & Less mature tooling \\
\hline
\end{tabular}
\end{center}

\subsection{Memory Cost of Optimizer State}

AdamW stores $m_t$ and $v_t$ at full precision for every parameter, so
optimizer state is roughly $2\times$ the size of the parameters (or
$\sim 8$ bytes/param at fp32).

\Ident\ Write the per-worker training memory budget, under A2.3, as
the sum of four terms:
\begin{equation}
\label{eq:ch2-mem-budget}
\boxed{\;M_{\mathrm{total}}
  \;=\; \underbrace{b_{\mathrm p}\, N}_{\text{params}}
     + \underbrace{b_{\mathrm g}\, N}_{\text{grads}}
     + \underbrace{b_{\mathrm o}\, N}_{\text{optimizer}}
     + \underbrace{M_{\mathrm{act}}(B,\,S,\,L,\,d_{\mathrm{model}})}_{\text{activations}}\;}
\end{equation}
with $b_{\mathrm p}, b_{\mathrm g}, b_{\mathrm o}$ bytes-per-parameter
coefficients for parameters, gradients, and optimizer state
respectively, and $M_{\mathrm{act}}$ the activation memory, which for
a Transformer scales roughly as $M_{\mathrm{act}} \propto
B\,S\,L\,d_{\mathrm{model}}$ with batch $B$, sequence length $S$,
layer count $L$, and model dimension $d_{\mathrm{model}}$.

\Emp\ For mixed-precision AdamW with an fp32 master copy, typical
coefficients are $b_{\mathrm p} = 2$ (bf16 weights in compute) $+ 4$
(fp32 master) $= 6$, $b_{\mathrm g} = 2$, $b_{\mathrm o} = 8$ (fp32
$m_t$ and $v_t$), giving
\[
M_{\mathrm{total}} \;\approx\; (6 + 2 + 8)\,N + M_{\mathrm{act}}
  \;=\; 16\,N + M_{\mathrm{act}}\;\text{bytes}.
\]
A $10^{10}$-parameter model therefore incurs $\sim 160$~GB of
resident state per worker \emph{before} activations — already beyond
a single accelerator's HBM.

\Cons\ This arithmetic motivates every memory-saving mechanism
covered later: ZeRO~\citep{rajbhandari2020zero} shards the $16N$
static term across workers; activation checkpointing attacks
$M_{\mathrm{act}}$; QLoRA and 8-bit optimizers~\citep{dettmers20228bit}
reduce $b_{\mathrm o}$; quantization at serving time (Chapter~8) reduces
$b_{\mathrm p}$ further. Chapter~7 revisits the sharded form of
Eq.~\eqref{eq:ch2-mem-budget}.

\section{Learning Rate Schedules}

The learning rate $\eta_t$ is the single most important hyperparameter
for Transformer training. Three families dominate LLM practice.

\subsection{Linear Warmup}

Start small and ramp up linearly over $T_w$ steps:
\begin{equation}
\eta_t \;=\; \eta_{\max} \cdot \min\!\left(1, \frac{t}{T_w}\right).
\end{equation}
Typical settings use $T_w \in [500, 5000]$ steps for models of
$10^{9}$--$10^{11}$ parameters. Warmup is essential because early
gradients are dominated by unstructured initialization: a full-size step
can drive parameters into a bad basin from which Adam's second-moment
estimator cannot recover quickly.

\subsection{Cosine Decay}

After warmup, the rate decays smoothly toward a minimum $\eta_{\min}$:
\begin{equation}
\eta_t \;=\; \eta_{\min} + \tfrac{1}{2}(\eta_{\max} - \eta_{\min})
            \left[1 + \cos\!\left(\pi\,\frac{t - T_w}{T_{\max} - T_w}\right)\right],
\qquad t > T_w.
\end{equation}
Cosine decay avoids the abrupt ``cliffs'' of step schedules and empirically
matches or exceeds more elaborate alternatives on LLM
pretraining~\citep{loshchilov2017sgdr}. A common variant holds
$\eta_{\min} = 0.1\,\eta_{\max}$ so that training never truly stops making
progress.

\subsection{Inverse-Square-Root (Noam) Schedule}

The original Transformer~\citep{vaswani2017attention} used
\begin{equation}
\eta_t \;=\; d_{\mathrm{model}}^{-1/2} \cdot
            \min\!\left(t^{-1/2},\; t\, T_w^{-3/2}\right),
\end{equation}
which combines linear warmup with an inverse-square-root decay.
Although cosine has largely displaced it for pretraining, inverse-sqrt
remains the default in many fine-tuning recipes.

\subsection{Engineering Guidance}

Choose schedules by matching the workload:
pretraining runs that consume a fixed compute budget pair warmup+cosine;
continual pretraining and fine-tuning that resume from a checkpoint
usually want warmup+constant or warmup+linear; very long-horizon runs
sometimes adopt \emph{one-cycle} / super-convergence
schedules~\citep{smith2018superconvergence} where $\eta$ rises and falls
once over the entire run.

Figure~\ref{fig:lr_schedule} overlays the three schedules discussed
above on a common axis; for the figure, recall from
\S\ref{sec:ch2-notation} the four schedule parameters:
warmup length $T_w$, total training horizon $T_{\max}$, peak rate
$\eta_{\max}$, and floor rate $\eta_{\min}$. Cosine decay uses all
four; inverse-square-root uses only $T_w$ and $\eta_{\max}$; the
constant baseline uses only $\eta_{\max}$.

\begin{figure}[H]
  \centering
  \begin{tikzpicture}[
  scale=1.0,
  axis/.style={->, thick, >=Stealth},
  curve/.style={thick, blue!60!black},
  warm/.style={thick, orange!80!black},
  cos/.style={thick, teal!60!black},
  const/.style={thick, red!60!black, dashed},
  lbl/.style={font=\footnotesize},
  phase/.style={font=\scriptsize\itshape}
]

% Axes
\draw[axis] (0,0) -- (10.3,0) node[right, lbl] {step $t$};
\draw[axis] (0,0) -- (0,4.3) node[above, lbl] {learning rate $\eta_t$};

% Peak LR reference
\draw[dotted] (0,3.5) -- (10,3.5);
\node[lbl, left] at (0,3.5) {$\eta_{\max}$};

% Warmup segment (linear up to t = 1.5)
\draw[warm] (0,0) -- (1.5, 3.5);
\node[phase, above, orange!60!black] at (0.75, 1.8) {warmup};

% Cosine decay from 1.5 to 9.5
\draw[cos] plot[domain=1.5:9.5, samples=70]
  (\x, { 0.3 + 0.5 * (3.5 - 0.3) * (1 + cos(((\x-1.5)/8)*180)) });
\node[phase, teal!60!black] at (5.5, 2.6) {cosine decay};

% Final constant tail
\draw[cos] (9.5,0.3) -- (10,0.3);

% Phase markers
\draw[dashed] (1.5,0) -- (1.5,3.5);
\node[phase, below] at (1.5,-0.1) {$t_w$};
\draw[dashed] (9.5,0) -- (9.5,0.3);
\node[phase, below] at (9.5,-0.1) {$T$};

% Constant LR (naive) overlay
\draw[const] (0,3.5) -- (10,3.5);
\node[lbl, right, red!60!black] at (7.2,3.8) {constant (naive)};

% Inverse-sqrt (Noam) overlay
\draw[thick, magenta!70!black, dotted] plot[domain=0.1:10, samples=60]
  (\x, {3.5 * min(\x/1.5, 1/sqrt(\x/1.5))});
\node[lbl, magenta!70!black, right] at (6.8, 1.5) {inverse-sqrt (Noam)};

\end{tikzpicture}

  \caption{Learning rate schedules. Linear warmup to $\eta_{\max}$ at
           step $T_w$, then (teal) cosine decay to floor $\eta_{\min}$
           at $T_{\max}$, or (magenta, dotted) inverse-square-root
           decay (Noam). The red dashed constant schedule is shown
           for contrast; it is the naive default that often diverges
           during early training on Transformer workloads.}
  \label{fig:lr_schedule}
\end{figure}

\section{Gradient Clipping}

Deep Transformer stacks can produce exploding-norm gradients when
activations or attention logits grow during early training. Gradient
clipping bounds the global step size:
\begin{equation}
\label{eq:ch2-gradclip}
g_t \;\leftarrow\; g_t \cdot \min\!\left(1, \frac{c}{\|g_t\|_2}\right),
\qquad c \in [0.5, 2.0].
\end{equation}
With the $\min$, the update direction is preserved while the magnitude
is capped at $c$. This is strictly preferable to elementwise clipping,
which distorts direction. Empirically, $c = 1.0$ is the default for
large LLM pretraining. Clipping is a cheap, principled defense against
rare but catastrophic optimizer divergence~\citep{pascanu2013difficulty},
and its activation rate is a useful operational telemetry signal: a
sudden increase typically precedes a loss spike by a few hundred steps.

\begin{workedexample}[clipping preserves direction, bounds magnitude]
Consider a two-parameter toy with raw gradient
$g_t = (3, 4) \in \R^{2}$, so $\|g_t\|_2 = 5$, and clip threshold
$c = 1$. Since $\|g_t\|_2 > c$, the scaling factor in
Eq.~\eqref{eq:ch2-gradclip} equals $c / \|g_t\|_2 = 1/5$, and the
clipped gradient is
\[
g_t^{\mathrm{clip}}
  \;=\; g_t \cdot \frac{1}{5}
  \;=\; \left(\tfrac{3}{5},\,\tfrac{4}{5}\right),
\qquad
\|g_t^{\mathrm{clip}}\|_2 \;=\; 1 \;=\; c.
\]
The unit vector $\hat g = g_t / \|g_t\|_2 = (0.6, 0.8)$ is \emph{exactly
preserved} since $g_t^{\mathrm{clip}} = c\,\hat g$; only the magnitude
changed. Contrast with elementwise clipping at the same threshold,
which would produce $g_t^{\mathrm{elem}} = (\min(3,1), \min(4,1)) = (1,1)$,
pointing along a direction $45^{\circ}$ off from $\hat g$. For a
second case, $g_t = (0.3, 0.4)$ has $\|g_t\|_2 = 0.5 < c$; the clip
factor is $\min(1, c/0.5) = 1$ and the gradient passes through
unchanged. This is why the activation rate of Eq.~\eqref{eq:ch2-gradclip}
(fraction of steps for which $\|g_t\|_2 > c$) is a clean telemetry
signal: under stable training it is stably small; under incipient
divergence it rises sharply well before loss explodes.
\end{workedexample}

\section{Mixed Precision, Loss Scaling, and Accumulation}

\subsection{Mixed-Precision Arithmetic}

LLM training uses mixed precision to increase throughput and reduce
memory~\citep{micikevicius2018mixed}: activations and gradients are held
in 16-bit (fp16 or bf16), while a ``master copy'' of parameters and
optimizer state is held in fp32. The chain is
\[
\theta_{\mathrm{fp32}} \xrightarrow{\text{cast}} \theta_{\mathrm{fp16}}
\xrightarrow{\text{forward/backward}} g_{\mathrm{fp16}}
\xrightarrow{\text{cast}} g_{\mathrm{fp32}}
\xrightarrow{\text{optimizer}} \theta_{\mathrm{fp32}}.
\]
The optimizer step is performed in fp32 to avoid accumulation error on
$m_t$, $v_t$. Bf16 is preferred over fp16 on modern hardware because its
wider exponent eliminates the need for loss scaling.

\subsection{Loss Scaling}

When using fp16 (not bf16), small gradients underflow to zero. Loss
scaling multiplies the loss by $s \in \{2^{10}, 2^{24}\}$ before
backpropagation, then divides gradients by $s$ before the optimizer
step. \emph{Dynamic} loss scaling halves $s$ whenever an $\mathrm{NaN}$
or $\mathrm{Inf}$ is detected and doubles it periodically otherwise.
This is an instance of adaptive numerical conditioning at the systems
layer, invisible to the mathematical model.

\subsection{Gradient Accumulation}

When the desired effective batch size exceeds accelerator memory,
gradients from $K$ micro-batches are summed in place before a single
optimizer step:
\[
g_t \;=\; \frac{1}{K}\sum_{k=1}^{K} g_t^{(k)}.
\]
Accumulation is algebraically equivalent to a $K$-times-larger batch at
the cost of $K$ extra forward/backward passes per update. It is the
primary mechanism by which organizations match the effective batch
sizes in published recipes while running on smaller clusters.

\subsection{Activation Checkpointing}

Activation memory dominates GPU usage. Activation checkpointing
\citep{chen2016training} trades compute for memory by discarding
intermediate activations during the forward pass and recomputing them
during backprop, typically saving $\sqrt{L}$ activation memory for $L$
layers at a $\sim 30\%$ compute overhead. This is how multi-hundred-billion
parameter models fit on realistic hardware.

\section{Overparameterization, Double Descent, and Benign Overfitting}

Classical statistical learning
theory~\citep{vapnik1998statistical,bishop2006prml} suggests that models
with more parameters than data should overfit catastrophically. Modern
deep networks contradict this: with $d \gg N$, training loss goes to
zero and test loss \emph{continues to improve}.

\subsection{Double Descent}

\Emp\ Plotted as a function of model size, test error exhibits the
classical U-curve up to the \emph{interpolation threshold} where
training error reaches zero. Beyond that threshold, test error
descends again~\citep{belkin2019reconciling,nakkiran2020double}. This
is an \emph{observed} behavior on standard image and language
benchmarks; the underlying theory connecting it to overparameterization,
implicit bias of SGD, and feature learning remains active research.
Three regimes:
\begin{itemize}
  \item \textbf{Underparameterized}: bias dominates; more capacity helps.
  \item \textbf{Interpolation threshold}: variance is maximized;
        test error spikes.
  \item \textbf{Overparameterized}: implicit regularization dominates;
        more capacity continues to help.
\end{itemize}
LLMs live deep in the third regime, which is why ``just train a bigger
model'' has been such a reliable strategy.

\subsection{Why Bigger Is Not Worse}

\Heur\ Several mechanisms are \emph{believed} to contribute to benign
overfitting; each is an engineering heuristic supported by partial
theory, not a proof:
\begin{itemize}
  \item SGD has an implicit bias toward low-norm, flat
        minima~\citep{keskar2017large}, which correlate with
        generalization;
  \item high-dimensional parameter spaces admit many zero-training-loss
        solutions, of which the optimizer tends to pick well-behaved
        ones;
  \item training data is structured and redundant, so capacity is used
        for coverage rather than memorization.
\end{itemize}
None of these arguments is a full theory, but together they justify the
engineering decision to scale up rather than regularize down.

\section{Scaling Laws}

\subsection{Kaplan Scaling}

\citet{kaplan2020scaling} observed that cross-entropy loss follows
smooth power laws in model size $N$, dataset size $D$, and compute $C$:
\begin{equation}
\label{eq:ch2-kaplan}
L(N) \;\approx\; \left(\frac{N_c}{N}\right)^{\alpha_N},
\qquad
L(D) \;\approx\; \left(\frac{D_c}{D}\right)^{\alpha_D},
\qquad
L(C) \;\approx\; \left(\frac{C_c}{C}\right)^{\alpha_C},
\end{equation}
with fitted exponents $\alpha_N \approx 0.076$,
$\alpha_D \approx 0.095$, and $\alpha_C \approx 0.05$. Equivalently,
\[
\log L \;=\; \mathrm{const} - \alpha_N \log N,
\]
so on log--log axes loss is a straight line in $N$ over six orders of
magnitude \emph{within the fitted compute regime}. This in-regime
predictability is what gives LLM pretraining its engineering
character; extrapolation outside the fitted compute range remains
a research question, not a guarantee.

\subsection{Chinchilla Scaling}

\citet{hoffmann2022chinchilla} refitted the joint law with better
hyperparameter tuning and reached a very different conclusion. Their
form is
\begin{equation}
\label{eq:ch2-chinchilla}
L(N, D) \;=\; E + \frac{A}{N^{\alpha}} + \frac{B}{D^{\beta}},
\end{equation}
with approximately $\alpha \approx 0.34$, $\beta \approx 0.28$,
$E \approx 1.69$ under A2.2.

\begin{derivation}[for the compute-optimal $(N^\star, D^\star)$]
Fix a compute budget $C$ with the Chinchilla approximation
$C \approx 6\,N D$ FLOPs (two FLOPs per MAC, three MACs per parameter
per token, forward plus backward). Minimize
\[
L(N, D) \;=\; E + \frac{A}{N^{\alpha}} + \frac{B}{D^{\beta}}
\qquad\text{subject to}\qquad
N D \;=\; \tfrac{C}{6}.
\]
Form the Lagrangian
$\mathcal{J}(N, D, \mu) = L(N, D) - \mu\,(N D - C/6)$. Stationarity
gives
\[
\partial_N \mathcal{J} \;=\; -\alpha\,A\,N^{-\alpha-1} - \mu\,D
\;=\; 0,
\qquad
\partial_D \mathcal{J} \;=\; -\beta\,B\,D^{-\beta-1} - \mu\,N
\;=\; 0.
\]
Divide the two first-order conditions to eliminate $\mu$:
\[
\frac{\alpha\,A\,N^{-\alpha-1}}{D} \;=\; \frac{\beta\,B\,D^{-\beta-1}}{N}
\qquad\Longleftrightarrow\qquad
\alpha\,A\,N^{-\alpha} \;=\; \beta\,B\,D^{-\beta},
\]
which fixes the \emph{ratio} between the two loss terms at the
optimum: they are matched up to the constant
$(\alpha A)/(\beta B)$. Solving for $D$ in terms of $N$ gives
$D^\star = \big[(\beta B)/(\alpha A)\big]^{1/\beta}\,N^{\alpha/\beta}$,
and substituting back into $N D = C/6$ yields
\[
N^\star \;\propto\; C^{a},
\qquad
D^\star \;\propto\; C^{b},
\qquad
a \;=\; \frac{\beta}{\alpha + \beta},
\qquad
b \;=\; \frac{\alpha}{\alpha + \beta}.
\]
With Chinchilla's fitted $\alpha \approx 0.34$, $\beta \approx 0.28$,
we get $a \approx 0.45$, $b \approx 0.55$; the ``equal exponents''
approximation $a \approx b \approx 0.5$ of
Eq.~\eqref{eq:ch2-chinchilla-opt} reflects the near-equality of
$\alpha$ and $\beta$ in the original fit.
\end{derivation}

\begin{equation}
\label{eq:ch2-chinchilla-opt}
N^\star \propto C^{a},
\qquad
D^\star \propto C^{b},
\qquad a \approx b \approx 0.5,
\end{equation}
i.e.\ model parameters and tokens should scale roughly equally with
compute. The Kaplan recommendation of ``scale model size faster than
dataset'' was an artifact of undertuned baselines. The practical
implication is dramatic: for any fixed compute budget there is a
\emph{cost-optimal} model size, and training a larger model on fewer
tokens is strictly worse than training a smaller model on more tokens.
Figure~\ref{fig:scaling_laws} shows the geometry: each model size
$N_i$ traces a loss curve in compute, and the compute-optimal
frontier is the envelope of the minima of those curves.

\begin{figure}[H]
  \centering
  \begin{tikzpicture}[
  scale=1.0,
  axis/.style={->, thick, >=Stealth},
  lbl/.style={font=\footnotesize},
  curv/.style={thick, smooth},
  frontier/.style={very thick, red!70!black, dashed}
]

% Axes (log-log conceptual; no numerical ticks)
\draw[axis] (0,0) -- (9.8,0) node[right, lbl] {compute $C$ (log)};
\draw[axis] (0,0) -- (0,5.3) node[above, lbl] {test loss $L$ (log)};

% Three "iso-model-size" curves: small, medium, large N
% Each is a decreasing curve that plateaus when N saturates.
% Small model N_1: plateau early
\draw[curv, blue!50!black] plot[domain=0.2:9, samples=70]
  (\x, {2.6 + 1.6 * exp(-0.45*\x) + 0.0001*\x*\x});
\node[lbl, blue!50!black] at (8.5, 2.7) {$N_1$ (small)};

% Medium N_2: plateau later
\draw[curv, orange!70!black] plot[domain=0.4:9, samples=70]
  (\x, {1.9 + 1.9 * exp(-0.32*\x)});
\node[lbl, orange!70!black] at (8.5, 2.0) {$N_2$};

% Large N_3: plateaus even later
\draw[curv, teal!50!black] plot[domain=0.8:9, samples=70]
  (\x, {1.3 + 2.0 * exp(-0.22*\x)});
\node[lbl, teal!50!black] at (8.5, 1.35) {$N_3$};

% Even larger N_4
\draw[curv, purple!60!black] plot[domain=1.2:9, samples=70]
  (\x, {0.8 + 2.2 * exp(-0.15*\x)});
\node[lbl, purple!60!black] at (8.5, 0.9) {$N_4$ (large)};

% Compute-optimal frontier (envelope of minima)
\draw[frontier] plot[domain=0.5:8.8, samples=50]
  (\x, {2.5 * exp(-0.20*\x) + 0.55});

\node[lbl, red!70!black] at (5.2, 3.2) {compute-optimal frontier};
\node[lbl, red!70!black] at (5.2, 2.85) {\textit{(Chinchilla: $N \propto C^{0.5}$, $D \propto C^{0.5}$)}};

% Label region above frontier
\node[lbl, gray] at (2.2, 4.4) {over-sized model, under-trained};
\node[lbl, gray] at (7.2, 4.0) {(wastes compute)};

% Label region below frontier
\node[lbl, gray] at (2.5, 0.4) {under-sized, over-trained};

\end{tikzpicture}

  \caption{Conceptual scaling-law picture. Each coloured curve is
           one model size trained on increasing compute; the
           dashed red envelope is the compute-optimal frontier
           along which Chinchilla balances $N$ and $D$. Regions
           above the envelope (too-big model, under-trained) and
           below it (too-small model, over-trained) both waste
           compute relative to the frontier.}
  \label{fig:scaling_laws}
\end{figure}

\subsection{What Scaling Laws Do and Do Not Tell You}

\Norm\ Scaling laws predict \emph{training loss} as a function of
resources under A2.2; they do \emph{not} predict:
\begin{itemize}
  \item emergent capabilities (discussed in Chapter~6 via
        \citet{wei2022emergent} and \citet{schaeffer2023mirage});
  \item performance on specific downstream tasks;
  \item the effect of data quality, deduplication, or mixture ratios.
\end{itemize}
They are compute-planning tools, not capability forecasts. Any
document that uses Eq.~\eqref{eq:ch2-chinchilla-opt} to argue for a
capability level, rather than a loss level, is conflating these two
predictive regimes. The minimum bar for using the law as a capability
predictor is an empirical fit on the specific task being claimed; in
the absence of such a fit, capability claims drawn from the joint-loss
law alone should be treated as out-of-scope.

\begin{provedassumed}
\emph{Proved (within the fit range).}
Eq.~\eqref{eq:ch2-chinchilla} is a regression fit whose exponents
follow as closed-form stationarity conditions under Lagrangian
optimization (the derivation above); the compute-optimal exponents
$a = \beta/(\alpha+\beta)$, $b = \alpha/(\alpha+\beta)$ are
mathematical consequences of the functional form.

\emph{Assumed.}
(i) The parametric form of Eq.~\eqref{eq:ch2-chinchilla} — a power
law plus a constant floor — is an empirical fit, not a derivation
from first principles.
(ii) The exponents $(\alpha, \beta, E)$ extrapolate only within the
compute regime they were fit on; frontier training at
$\gg 10^{24}$~FLOPs is an \emph{extrapolation} whose validity is
itself an empirical question.
(iii) The FLOP-cost approximation $C \approx 6\,N D$ ignores
inference-time, activation-recompute, and MoE routing cost, and is
only a first-order estimate.
(iv) Loss is a proxy for capability; the gap between the two is
what Chapter~6 calls ``emergent capabilities'' and is \emph{not}
governed by Eq.~\eqref{eq:ch2-chinchilla}.
\end{provedassumed}

\section{Systems Engineering Implications}

\begin{center}
\begin{tabular}{|p{3.4cm}|p{3.4cm}|p{6.8cm}|}
\hline
V-model stage & Optimizer / training property & System requirement or activity \\
\hline
Requirements & Fixed compute budget & Use Chinchilla law to size $N, D$. \\
 & Reproducibility target & Seeded data order, deterministic kernels. \\
\hline
Architecture & AdamW state $\sim 2\times$ params & Budget optimizer memory (ZeRO). \\
 & fp16 vs bf16 & Choose based on hardware and stability. \\
 & Activation memory & Plan recompute vs sharding. \\
\hline
Implementation & Warmup + cosine & Expose schedule params to config. \\
 & Gradient accumulation & Log effective vs nominal batch. \\
\hline
V\&V & Loss spikes and plateaus & Divergence alarms; checkpoint rollback. \\
 & Clip activation rate & Telemetry feed into regression gates. \\
 & Small-scale proxies & Shadow runs at $10^{-3}\times$ compute
   \citep{wortsman2023stable}. \\
\hline
Operations & Regressions between runs & Rerun on fixed seed + data; bisect. \\
 & Cost / quality tracking & Dollars per perplexity point. \\
\hline
\end{tabular}
\end{center}

\section{Reproducibility in Distributed Training}

Deterministic training under data-parallel, tensor-parallel, and
pipeline-parallel execution is substantially harder than it sounds.
Common sources of nondeterminism:
\begin{itemize}
  \item asynchronous all-reduce summation order;
  \item nondeterministic cuDNN kernels (e.g., convolutions);
  \item dataloader worker scheduling;
  \item mixed-precision dynamic loss scaling decisions;
  \item floating-point non-associativity across restart boundaries.
\end{itemize}
A defensible reproducibility protocol fixes:
(i)~the random seed for parameter init, data shuffling, and dropout;
(ii)~the global data order (checkpoint the dataloader state);
(iii)~the kernel selection via deterministic-mode flags;
(iv)~the loss-scale trajectory (deterministic schedule rather than
dynamic halving).
Even then, \emph{bit-exact} reproducibility across hardware generations
is rarely achievable; the practical goal is \emph{statistical}
reproducibility, i.e., loss curves within noise bands across reruns.

\section{Human Factors and Failure Modes}

A recurring failure mode is the attribution of \emph{training behavior}
to \emph{model intelligence}. Loss plateaus, loss drops after restarts,
apparent ``insight moments'' — all of these are best explained as
optimization dynamics. Concrete examples:
\begin{itemize}
  \item Attributing a mid-training loss jump to a capability emerging,
        when the actual cause was an inadvertent change in data
        ordering after a checkpoint restore.
  \item Interpreting a flat training loss as ``the model has learned
        everything it can,'' when the actual cause is a learning-rate
        schedule that has decayed too far.
  \item Overreading the output of small-scale probes: noise at
        $10^9$ parameters does not monotonically forecast behavior at
        $10^{11}$ parameters.
\end{itemize}
The countermeasure is disciplined telemetry: track gradient norm,
clip-rate, parameter norm, activation statistics, and token-ordering
hash on every run. Most training ``mysteries'' resolve to an operator
mistake that better instrumentation would have caught.

\section{Bridges to Later Chapters}

\begin{itemize}
  \item \textbf{Chapter 4 (Tokenization and Embeddings)}: embedding tables
        are the largest single contributor to optimizer state memory;
        tokenizer choice directly affects the effective $D$ in
        Eq.~\eqref{eq:ch2-chinchilla}.
  \item \textbf{Chapter 5 (Transformer Architecture)}: the layer
        normalization placement (pre-LN vs post-LN)
        \citep{xiong2020onlayer} interacts with warmup schedule length.
  \item \textbf{Chapter 6 (Pretraining \& In-Context)}: the scaling-law
        formulas set the budget; emergent capability claims must be
        checked against them for confounds~\citep{schaeffer2023mirage}.
  \item \textbf{Chapter 7 (Efficiency, Scaling, MoE)}: ZeRO
        \citep{rajbhandari2020zero} and pipeline parallelism
        \citep{narayanan2021efficient} are direct responses to the
        optimizer-memory constraint derived here.
  \item \textbf{Chapter 8 (Adaptation, LoRA, Quantization)}: adaptation
        methods replace a full-fidelity optimizer over $10^{11}$
        parameters with one over $\le 10^{7}$, tracing back to the
        cost equation laid out in \S\,Scaling Laws.
  \item \textbf{Chapter 11 (Governance \& Assurance)}: reproducibility
        protocols become evidence in an assurance case; training
        regression detection is an operational-safety control.
\end{itemize}

\section{Key References}

\begin{itemize}
  \item \citet{goodfellow2016dl} — canonical treatment of optimization
        and generalization in deep learning.
  \item \citet{kingma2015adam} — original Adam paper.
  \item \citet{loshchilov2019decoupled} — AdamW, decoupled weight
        decay.
  \item \citet{loshchilov2017sgdr} — cosine-decay schedules.
  \item \citet{goyal2017accurate} — large-batch warmup and the linear
        scaling rule.
  \item \citet{pascanu2013difficulty} — origin of gradient clipping.
  \item \citet{micikevicius2018mixed} — mixed-precision training.
  \item \citet{chen2016training} — activation checkpointing.
  \item \citet{kaplan2020scaling} — scaling laws for neural language
        models.
  \item \citet{hoffmann2022chinchilla} — compute-optimal scaling.
  \item \citet{nakkiran2020double}, \citet{belkin2019reconciling} —
        double descent and benign overfitting.
  \item \citet{keskar2017large} — large-batch training and flat
        minima.
  \item \citet{zhang2017understanding} — empirical study of
        overparameterization.
  \item \citet{wortsman2023stable} — small-scale proxies for
        large-training instabilities.
\end{itemize}

\section{Recommended University Lectures}

\begin{itemize}
  \item MIT 6.S191 --- optimization and deep learning practice,
        \citet{mit_6s191}.
  \item Berkeley CS182 --- deep learning foundations,
        \citet{berkeley_cs182}.
  \item Stanford CS229 --- gradient descent and generalization,
        \citet{stanford_cs229}.
  \item Stanford CS324 --- Large Language Models, module on training
        dynamics, \citet{stanford_cs324}.
  \item Caltech CS/CNS/EE~156 --- Learning from Data,
        \citet{caltech_cs156}.
\end{itemize}

\section{Practitioner Lab}

\begin{enumerate}
  \item \textbf{Optimizer comparison.} Train a $\sim\!10$\,M-parameter
        Transformer on WikiText-2 with SGD, SGD+momentum, Adam, and
        AdamW. Plot training and validation loss; report steps-to-target
        and final perplexity.
  \item \textbf{Bias correction.} Remove the
        $1-\beta_1^{\,t}, 1-\beta_2^{\,t}$ terms from
        Eq.~\eqref{eq:ch2-adam-update}. Measure the impact on the first
        500 steps.
  \item \textbf{LR sweep.} Sweep $\eta_{\max} \in \{1, 3, 10, 30, 100\}
        \times 10^{-5}$ with a warmup+cosine schedule. Plot loss at step
        1{,}000 vs $\eta_{\max}$ to locate the instability knee.
  \item \textbf{Gradient clipping telemetry.} Log the clip activation
        rate and gradient norm $\|g_t\|_2$ every step. Correlate spikes
        in $\|g_t\|_2$ with subsequent loss jumps.
  \item \textbf{Mini-Chinchilla.} At fixed compute $C$, train
        $(N, D)$ pairs $(N, 20 N)$ and $(2N, 10 N)$ and $(N/2, 40 N)$.
        Verify that $(N, 20 N)$ yields the lowest loss.
  \item \textbf{Reproducibility audit.} Run the same training script
        twice with the same seed; diff the loss curves. Identify and
        eliminate all nondeterminism sources until the curves match to
        within numerical tolerance.
\end{enumerate}

\section*{Derivation Ledger (Ch. 2)}

\begin{derivledger}
Compact index of the chapter's central identities; each cites the full
derivation already given in the text.
\begin{itemize}
  \item \textbf{Empirical-risk objective.}\
        $\mathcal{L}(\theta) = \tfrac{1}{N_{\mathrm{tok}}} \sum_i \ell(\theta; x_i)$
        --- Eq.~\eqref{eq:ch2-empirical-loss}.
  \item \textbf{Minibatch SGD update.}\
        Eq.~\eqref{eq:ch2-sgd}; assumption (unbiased gradient
        estimate, bounded variance) stated in §SGD-to-SDE.
  \item \textbf{SGD-to-SDE (Langevin approximation).}\
        $d\theta_t = -\nabla \mathcal{L}\, dt + \sqrt{\eta \Sigma}\, dW_t$
        --- Eq.~\eqref{eq:ch2-langevin}; requires small learning rate
        and approximately Gaussian gradient noise.
  \item \textbf{Adam / AdamW updates.}\
        First and second moment recursions
        Eqs.~\eqref{eq:ch2-adam-m}--\eqref{eq:ch2-adam-update};
        decoupled weight-decay form Eq.~\eqref{eq:ch2-adamw-update}.
  \item \textbf{Mixed-precision memory budget.}\
        $M = (b_{\rm param} + b_{\rm grad} + b_{\rm opt}) \cdot N + M_{\rm act}$
        --- Eq.~\eqref{eq:ch2-mem-budget}; reused by CH7 for LoRA and
        QLoRA budgets.
  \item \textbf{Gradient clipping.}\
        Rescale $g \leftarrow g \cdot \min(1, \tau / \|g\|)$
        --- Eq.~\eqref{eq:ch2-gradclip}.
  \item \textbf{Kaplan scaling law.}\
        Power-law loss in $N, D, C$ --- Eq.~\eqref{eq:ch2-kaplan}.
  \item \textbf{Chinchilla compute-optimal.}\
        Lagrangian over $L(N,D)$ under $C = 6ND$ yields
        $N^* \propto C^a$, $D^* \propto C^{1-a}$
        --- Eqs.~\eqref{eq:ch2-chinchilla}--\eqref{eq:ch2-chinchilla-opt}.
\end{itemize}
\end{derivledger}

\section{Chapter 2 Takeaway}

\begin{quote}
Training behavior in LLMs is governed less by ``intelligence'' than by
optimization dynamics operating in very high dimensions. Adam/AdamW,
warmup, cosine decay, gradient clipping, and mixed precision are not
interchangeable tuning knobs; each responds to a specific pathology of
large-scale stochastic optimization, and scaling laws turn the resulting
process into a compute-planning exercise that can be reasoned about
engineering-first.
\end{quote}
